% ============================================================
% Оформление формул по ГОСТ
% Формулы выравниваются по левому краю с абзацным отступом.
% Пояснение символов оформляется окружением conditions.
% Подключать ПОСЛЕ \usepackage{amsmath}
% ============================================================

\usepackage{array}

% --- Вертикальный отступ текста от линий в ячейках таблиц ---
\setlength{\extrarowheight}{3pt}

% --- Отступ формул = абзацный отступ ---
\setlength{\mathindent}{\parindent}

% --- Окружение «где» для пояснения символов ---
% Использование:
%   \begin{conditions}
%     \sigma & нормальное напряжение, \textit{Па}; \\
%     N_z    & нормальная сила, \textit{Н}; \\
%     A      & площадь поперечного сечения стержня, \textit{м}$^2$.
%   \end{conditions}
\newenvironment{conditions}
  {\par\noindent где\quad
   \begin{tabular}[t]{>{$\displaystyle}l<{$} @{\ --\ } p{10.5cm}}}
  {\end{tabular}\par}
